% Diagrama de la arquitectura del modelo (figura de la sección «Primer modelo»).
% Cómo editarlo: los textos están en los \node de las etiquetas (al final de cada bloque); los colores, en
%   el bloque «Colores»; el tamaño de cada capa, en los argumentos de \cubo (ancho, alto, profundidad) y de \capavec.
% Cómo regenerarlo (desde esta carpeta):   pdflatex arquitectura.tex   -> arquitectura.pdf (el reporte lo incluye).
% Dependencias: TeX Live con tikz y standalone (los mismos paquetes del reporte).
\documentclass[tikz,border=3pt]{standalone}
\usepackage[spanish,es-noshorthands]{babel}
\usepackage[utf8]{inputenc}
\usepackage[T1]{fontenc}
\usepackage[dvipsnames]{xcolor}
\usepackage{amsmath}
\usetikzlibrary{calc,patterns,arrows.meta}

% Colores (los del reporte, más un acento para la salida)
\definecolor{skyblue}{RGB}{122,177,226}
\definecolor{daphne}{RGB}{51,118,185}
\definecolor{AirForceBlue}{RGB}{21,58,89}
\definecolor{midnight}{RGB}{18,34,62}
\definecolor{acento}{RGB}{231,138,43}
\definecolor{casillaclara}{RGB}{240,217,181}
\definecolor{casillaoscura}{RGB}{181,136,99}

% \cubo{x}{y}{ancho}{alto}{profundidad}{color}: volumen con la cara frontal en (x,y) y sus tres caras visibles
\newcommand{\cubo}[6]{%
  \pgfmathsetmacro{\dx}{0.55*#5}\pgfmathsetmacro{\dy}{0.40*#5}%
  \filldraw[fill=#6!55!white, draw=midnight, line width=0.3pt]
    (#1+#3,#2) -- (#1+#3+\dx,#2+\dy) -- (#1+#3+\dx,#2+#4+\dy) -- (#1+#3,#2+#4) -- cycle;
  \filldraw[fill=#6!30!white, draw=midnight, line width=0.3pt]
    (#1,#2+#4) -- (#1+\dx,#2+#4+\dy) -- (#1+#3+\dx,#2+#4+\dy) -- (#1+#3,#2+#4) -- cycle;
  \filldraw[fill=#6, draw=midnight, line width=0.3pt] (#1,#2) rectangle (#1+#3,#2+#4);}

% \capavec{x}{y}{alto}{color}: capa como barra vertical de neuronas
\newcommand{\capavec}[4]{%
  \filldraw[fill=#4, draw=midnight, line width=0.3pt, rounded corners=1.2pt] (#1,#2) rectangle (#1+0.34,#2+#3);}

\begin{document}
\begin{tikzpicture}[font=\sffamily\footnotesize, >={Stealth[length=4pt]}, flecha/.style={draw=midnight!70, line width=0.7pt, ->}]

  % --- Entrada: tablero de 8x8 casillas con 17 canales ---
  \cubo{0}{-0.8}{1.6}{1.6}{0.45}{skyblue}
  \foreach \i in {0,...,7}{\foreach \j in {0,...,7}{
    \pgfmathparse{mod(\i+\j,2)==0 ? 1 : 0}
    \ifnum\pgfmathresult=1 \fill[casillaoscura] (0.2*\i,-0.8+0.2*\j) rectangle +(0.2,0.2);
    \else \fill[casillaclara] (0.2*\i,-0.8+0.2*\j) rectangle +(0.2,0.2);\fi}}
  \draw[midnight, line width=0.3pt] (0,-0.8) rectangle (1.6,0.8);
  \node[align=center] at (1.0,-1.55) {\bfseries Tablero\\[-1pt]$8\times8\times17$};

  % --- Capas convolucionales ---
  \draw[flecha] (2.25,0) -- (3.0,0);
  \cubo{3.15}{-0.8}{1.6}{1.6}{1.1}{daphne}
  \draw[white!85!daphne, line width=0.25pt] foreach \k in {1,...,7}{(3.15+0.2*\k,-0.8) -- (3.15+0.2*\k,0.8) (3.15,-0.8+0.2*\k) -- (4.75,-0.8+0.2*\k)};
  \node[align=center] at (4.0,-1.55) {\bfseries Capa convolucional\\[-1pt]$3\times3$, 64 filtros\\[-1pt]$8\times8\times64$};

  \draw[flecha] (5.6,0) -- (6.3,0);
  \cubo{6.45}{-0.8}{1.6}{1.6}{1.8}{AirForceBlue}
  \draw[white!80!AirForceBlue, line width=0.25pt] foreach \k in {1,...,7}{(6.45+0.2*\k,-0.8) -- (6.45+0.2*\k,0.8) (6.45,-0.8+0.2*\k) -- (8.05,-0.8+0.2*\k)};
  \node[align=center] at (7.45,-1.55) {\bfseries Capa convolucional\\[-1pt]$3\times3$, 128 filtros\\[-1pt]$8\times8\times128$};

  % --- Aplanado y capas densas ---
  \draw[flecha] (9.2,0) -- node[above=1pt]{\itshape Flatten} (10.35,0);
  \capavec{10.5}{-1.5}{3.0}{skyblue}
  \node[align=center] at (10.67,-1.95) {\bfseries Vector\\[-1pt]8,192};

  \draw[flecha] (11.05,0) -- (11.85,0);
  \capavec{12.0}{-1.0}{2.0}{daphne}
  \node[align=center] at (12.17,-1.45) {\bfseries Densa\\[-1pt]2,048};

  % Dropout: neuronas apagadas (rayadas) durante el entrenamiento
  \draw[flecha] (12.7,0) -- (13.5,0);
  \foreach \y in {-1.0,-0.6,-0.2,0.2,0.6}{\filldraw[fill=daphne!80!white, draw=midnight, line width=0.3pt] (13.65,\y) rectangle +(0.34,0.4);}
  \foreach \y in {-0.6,0.2}{\filldraw[pattern=north east lines, pattern color=midnight!70, draw=midnight, line width=0.3pt] (13.65,\y) rectangle +(0.34,0.4);}
  \node[align=center] at (13.82,-1.45) {\itshape\bfseries Dropout\\[-1pt]0.5};

  \draw[flecha] (14.35,0) -- (15.15,0);
  \capavec{15.3}{-1.5}{3.0}{AirForceBlue}
  \node[align=center] at (15.47,-1.95) {\bfseries Densa\\[-1pt]4,096 (\textit{softmax})};

  % --- Salida: probabilidad de cada jugada ---
  \draw[flecha] (15.85,0) -- (16.35,0);
  \foreach \h/\y in {0.95/0.55,0.55/0.15,0.30/-0.25,0.15/-0.65}{
    \fill[acento] (16.45,\y) rectangle +(\h,0.3);}
  \node[align=center, text=acento!70!black] at (17.0,-1.2) {\bfseries Probabilidad de\\[-1pt]cada jugada};

\end{tikzpicture}
\end{document}
